%==============================================================
%  TÓM TẮT / ABSTRACT
%==============================================================

\chapter*{Tóm tắt}
\addcontentsline{toc}{chapter}{Tóm tắt}
\markboth{Tóm tắt}{Tóm tắt}

Luận văn trình bày một cách hệ thống cơ sở toán học, thuật toán và thực nghiệm
của \emph{mạng nơ-ron thông tin vật lý} (physics-informed neural network, PINN)
khi được dùng làm công cụ giải phương trình vi phân.

Điểm khác biệt về trọng tâm so với tài liệu học sâu tổng quát là luận văn khảo
sát mạng truyền thẳng từ góc độ \emph{đạo hàm theo biến đầu vào} chứ không phải
theo tham số. Hai hệ thức truy hồi cho Jacobi và Hessian
(Mệnh đề~\ref{md:jacobian}, \ref{md:hessian}) được thiết lập tường minh và trở
thành công cụ trung tâm cho mọi tính toán phần dư về sau; từ chúng suy ra ngay
một tiêu chuẩn \emph{cần} đối với hàm kích hoạt. Trường hợp ReLU được phát biểu
chính xác ở dạng mạnh hơn nhiều một sự ``suy giảm'': hàm mục tiêu phần dư của
phương trình bậc hai là \emph{hằng số địa phương}, nên mọi thuật toán dựa trên
gradient đứng yên tuyệt đối (Hệ quả~\ref{hq:pinn-hong}).

Về phương pháp, luận văn trả lời câu hỏi thường bị bỏ qua: \emph{vì sao cực tiểu
hoá phần dư lại là việc đáng làm}. Câu trả lời là một chặn ổn định của bài toán
vi phân (Mệnh đề~\ref{md:chan-on-dinh}), và hằng số của chặn ấy suy biến theo
$1/\nu$ với bài toán nhiễu kỳ dị (Mệnh đề~\ref{md:chan-suy-bien}). Kết hợp với
bệnh lý gradient ở tầng tối ưu và thiên kiến phổ ở tầng xấp xỉ, luận văn tổng
hợp \emph{ba nguyên nhân độc lập, nằm ở ba tầng khác nhau}, cùng làm PINNs suy
giảm trên đúng lớp bài toán có thang độ dài nhỏ (Bảng~\ref{tab:ba-nguyen-nhan}).

Phần thực nghiệm cài đặt PINNs bằng PyTorch, không dùng thư viện đóng gói sẵn,
và kiểm chứng từng phát biểu lý thuyết bằng bảy thí nghiệm có tiêu chí thành bại
xác định trước. Kết quả được đối chiếu với lý thuyết ở
Bảng~\ref{tab:doi-chieu}. Luận văn báo cáo đầy đủ cả những kết quả không đạt kỳ
vọng: một kết quả buộc phải siết lại cách phát biểu của chính chương lý thuyết
(Nhận xét~\ref{nx:hieu-chinh}), và một kết quả tiêu cực về thuật toán cân bằng
trọng số đã dẫn tới một chứng minh mới về chế độ phản hồi dương của quy tắc ấy
(Bổ đề~\ref{bd:phan-hoi-duong}).

\bigskip
\noindent\textbf{Từ khoá:} mạng nơ-ron thông tin vật lý, phương trình vi phân,
vi phân tự động, xấp xỉ phổ quát, thiên kiến phổ, bệnh lý gradient, chặn ổn
định, L-BFGS.

%--------------------------------------------------------------
\chapter*{Abstract}
\addcontentsline{toc}{chapter}{Abstract}
\markboth{Abstract}{Abstract}

\begin{english}
This thesis gives a systematic account of the mathematical, algorithmic and
experimental foundations of \emph{physics-informed neural networks} (PINNs) used
as solvers for differential equations.

Its focus differs from that of general deep-learning texts: feedforward networks
are studied through their derivatives \emph{with respect to the input} rather
than to the parameters. Two explicit recurrences for the input Jacobian and
Hessian are established and used throughout for residual computations; from them
follows a \emph{necessary} condition on the activation function. The ReLU case is
stated in a form considerably stronger than mere degradation: for a second-order
operator the residual objective is \emph{locally constant}, so every
gradient-based algorithm is exactly stationary.

Methodologically, the thesis answers a question that is usually skipped: \emph{why
minimising the residual is worth doing at all}. The answer is a stability bound
for the underlying differential problem, whose constant degenerates as $1/\nu$ for
singularly perturbed problems. Combined with gradient pathology at the
optimisation level and spectral bias at the approximation level, this yields
\emph{three independent causes, situated at three different levels}, that jointly
degrade PINNs on precisely the class of problems with small length scales.

The experimental part implements PINNs in PyTorch without any packaged PINN
library and tests each theoretical claim through six experiments with success
criteria fixed in advance. Negative results are reported in full: one experiment
forced a restatement of a conclusion in the theoretical chapter, and a negative
result on adaptive loss balancing led to a new proof that the balancing rule has
a positive-feedback failure mode.

\bigskip
\noindent\textbf{Keywords:} physics-informed neural networks, differential
equations, automatic differentiation, universal approximation, spectral bias,
gradient pathology, stability bound, L-BFGS.
\end{english}
